\section{Geometric origins and the calibrated heat realization}
\label{sec:geometry}

The preceding source can be described by geometric traces, scalar
potentials and a prepared operator orbit. These descriptions explain the
origin of the investigation. We derive the identities needed for the
Green thread while retaining their preparation and observation assumptions
\cite{GNPBackground,OperatorProof,OperatorAssessment}.

\subsection{The cylinder and the multiplicity-cancelling trace}
Let \(D\) be the Dirac operator on the unit two-sphere. Its eigenvalues
are \(\pm n\), each of multiplicity \(2n\), \(n\ge1\)
\cite{CamporesiHiguchi1996}. On the cylinder
\(\R_v\times S^2_{1/2}\), with the radial Clifford matrix anticommuting
with the angular operator, the geometric square is
\begin{equation}\label{eq:cylinder-square}
 L_0=-\partial_v^2+4A,\qquad A=D^2.
\end{equation}
The scalar curvature is \(8\), so the Lichnerowicz curvature term is
\(2\); it is already contained in this geometric square. The
\(A=n^2\) shell has total rank \(4n\). Its radial Green kernel is
\[
 g_n(v,w)=\frac{\ee^{-2n|v-w|}}{4n}.
\]
Indeed it solves \((-\partial_v^2+4n^2)g_n=\delta_w\): the derivative
jump at \(w\) is \(-1\). Thus the coincident kernel induces the angular
weight \(W=(4\sqrt A)^{-1}\), and
\begin{equation}\label{eq:geometric-trace}
 Q(\beta)=\Tr(W\ee^{-\beta A})=\sum_{n\ge1}\ee^{-\beta n^2}.
\end{equation}
All traces at \(\beta>0\) converge absolutely. The factor \(4n\) from
the shell cancels the diagonal Green factor \(1/(4n)\).

Supply the scan \(\beta(u)=\pi\ee^{4u}\). Then
\(p(u)=\ee^u Q(\beta(u))\), and differentiating twice gives
\begin{equation}\label{eq:trace-source}
 \Phi(u)=\ee^u Q(\beta)
 \Tr\bigl[\rho_\beta(2\beta^2A^2-3\beta A)\bigr],\qquad
 \rho_\beta=\frac{W\ee^{-\beta A}}{Q(\beta)}.
\end{equation}
The differentiation is justified by Gaussian angular summability on
compact \(u\) sets. This proves the positive-half-line source formula
and, together with theta reflection, the full source. The scan is a
specified relation between variables: translation of the radial Green
probe does not generate it, since \(g_n(v,v)\) is constant in \(v\).

\subsection{A scalar potential and its transported source}
Write \(\Theta(R)=1+2Q(\pi R^2)\). Jacobi reflection says
\(\Theta(R)=R^{-1}\Theta(1/R)\). The scalar potential
\(V(R)=1+R^{-1}-\Theta(R)\) gives the even function
\begin{equation}\label{eq:scalar-G0}
 \mathscr G_0(u)=\tfrac12\ee^u V(\ee^{2u})
 =\tfrac12\ee^{-u}-p(u)\quad(u\ge0).
\end{equation}
Let \(L=1-\partial_u^2\). The reflection endpoint \(p'(0)=-1/2\)
makes the derivative of \eqref{eq:scalar-G0} zero at the origin, so its
even extension has no distributional delta there. Hence
\(L\mathscr G_0=8\Phi\). More generally define
\begin{equation}\label{eq:scalar-G-lambda}
 \mathscr G_\lambda=8L^{-1}W_\lambda
 =4\int_\R \ee^{-|u-w|}W_\lambda(w)\dd w.
\end{equation}
The positive convolution and the complete source bounds justify every
fixed derivative and yield
\(\mathscr G_\lambda(u)\sim c_\lambda\ee^{-|u|}\), where
\(c_\lambda=8H_\lambda(\ii)>0\) and \(c_0=1/2\).
For example, as \(u\to+\infty\), the part \(w\le u\) equals
\(4\ee^{-u}\int_{-\infty}^u\ee^w W_\lambda(w)\dd w\); the complementary
tail is negligible by double-exponential decay. Evenness proves the
other end.

Fourier transformation of \(L\mathscr G_\lambda=8W_\lambda\) gives
\begin{equation}\label{eq:scalar-fourier}
 16H_\lambda(z)=(1+z^2)\int_\R\mathscr G_\lambda(u)\ee^{\ii zu}\dd u,
 \qquad |\Im z|<1.
\end{equation}
The raw integral is only asserted in this strip. At \(z=\pm\ii\) one
must use the entire source definition of \(H_\lambda\), rather than
multiply a vanishing polynomial by a divergent integral.

If \(X\) denotes multiplication by \(u\), source transport induces
\begin{equation}\label{eq:scalar-transport}
 \partial_\lambda\mathscr G_\lambda=\mathcal B\mathscr G_\lambda,
 \quad \mathcal B=L^{-1}X^2L,
 \quad \mathcal Bf=u^2f+L^{-1}(4uf'+2f).
\end{equation}
The commutator identity follows by expanding \(L(u^2f)\).
On Schwartz functions \(\mathcal B^j=L^{-1}X^{2j}L\); the displayed
prepared family supplies its own domains through \eqref{eq:scalar-G-lambda}.
Multiplying \(\mathscr G_0\) by \(\ee^{\lambda u^2}\) is a different
operation and is not integrable for \(\lambda>0\). This distinction is
essential when interpreting a geometric potential as a Newman source.

\subsection{A fixed operator with the correct heat clock}
The calibrated operator is
\begin{equation}\label{eq:calibrated-K}
 K=L_0-4I=-\partial_v^2+4(A-I).
\end{equation}
Under the unitary radial Fourier transform its shell multiplier is
\(m_n(p)=p^2+4(n^2-1)\), with maximal domain
\[
 \Dom K=\left\{\psi:\sum_n\int_\R m_n(p)^2
             \|\widehat\psi_n(p)\|^2\dd p<\infty\right\}.
\]
Real nonnegative multiplication proves self-adjointness and positivity;
the spectrum is purely absolutely continuous and equals \([0,\infty)\).
The value zero is not an eigenvalue. The semigroup \(S_\tau=\ee^{-\tau K}\)
has norm one, and for \(j\ge1,\tau>0\),
\(\|K^jS_\tau\|\le(j/(\ee\tau))^j\), by maximizing
\(r^j\ee^{-\tau r}\) over \(r\ge0\).

Fix a unit vector \(\chi\) in the first angular shell. Define
\(\widehat{E_\chi f}(p)=f(p)\chi\) and
\(R_\chi\psi(x)=\sqrt{\pi/2}\langle\chi,\psi(x)\rangle\).
The first is an isometry and the second is bounded between the relevant
\(L^2\) spaces. On their natural domains,
\begin{equation}\label{eq:operator-intertwine}
 KE_\chi=E_\chi X^2,\qquad R_\chi K=-\partial_x^2R_\chi.
\end{equation}
The complete prepared amplitude \(W_{\lambda_0}\) therefore satisfies
\begin{equation}\label{eq:operator-response}
 R_\chi S_\tau E_\chi W_{\lambda_0}=H_{\lambda_0-\tau}=h(\tau,\cdot).
\end{equation}
To verify the constant, inverse Fourier transformation gives
\(\sqrt{\pi/2}/\sqrt{2\pi}=1/2\), exactly the coefficient in
\eqref{eq:H}. Evenness removes the Fourier-sign convention. The source
bounds place the prepared orbit in every fixed graph domain. Its complex
frequency continuation and pointwise derivatives follow from weighted
source integrals, not bounded point evaluation on arbitrary \(L^2\).

Using \(L_0\) instead would produce \(\ee^{-4\tau}h\). The calibration
is material. In fact a fixed bounded linear observation of
\(\ee^{-\tau L_0}\psi\), with fixed \(\psi\), decays exponentially in
\(L^2\), whereas
\[
 \|H_{\lambda_0-\tau}\|_2
 \sim\Phi(0)(\pi/2)^{3/4}\tau^{-1/4}.
\]
This follows from Plancherel and scaling
\(\int\ee^{-2\tau u^2}|W_{\lambda_0}(u)|^2\dd u\).
It excludes this particular uncalibrated bounded-observation realization
for all \(\tau\), not all alternative preparations or observations.
Likewise reverse heat is unbounded on ambient states, while the
canonical source itself is entire in \(\lambda\).

The representation preserves its source-accuracy burden exactly:
\begin{equation}\label{eq:operator-error}
 \|R_\chi S_\tau E_\chi(f-g)\|_2
 =\sqrt{\pi/2}\,\|\ee^{-\tau u^2}(f-g)\|_2.
\end{equation}
No arithmetic cancellation estimate is supplied merely by changing
coordinates to this operator realization.

\subsection{Energy remains regular at a collision}
For the prepared orbit put
\(N=\int|W_{\lambda_0-\tau}|^2=(2/\pi)\int h^2\) and
\(T=\int u^2|W_{\lambda_0-\tau}|^2=(2/\pi)\int h_x^2\).
Differentiation and Plancherel give
\begin{equation}\label{eq:energy-laws}
 N'=-2T,\qquad T'=-2\|K S_\tau E_\chi W_{\lambda_0}\|_2^2,
 \qquad (T/N)'=-2\operatorname{Var}_{|W|^2/N}(u^2).
\end{equation}
The linear source probability \(W/\int W\) and the spectral probability
\(|W|^2/N\) are distinct; the latter produces the factor two above.
Locally, \((h^2)_\tau=(h^2)_{xx}-2h_x^2\). At a double zero with
\(h=h_x=0\), \(h_{xx}=a\ne0\), one has
\((h^2)_\tau=0\) and \((h^2)_{\tau\tau}=2a^2\).
The regular energy laws are compatible with a change in the real zero
configuration and impose no collision barrier.

\subsection{The other parent realizations and their scope}
For completeness, the shared parent also supplied two exact observation
models \cite{TNDDirect,TNDAux}. They are not new inputs to a Green sign
estimate. In the direct model prepare
\(f_\lambda=\sqrt{W_\lambda/Z_\lambda}\), \(Z_\lambda=2H_\lambda(0)\),
and supply the measured quadrature. With
\(U(t)=\ee^{-\ii\omega tL_0}\), \(Q_t=U(t)XU(t)^*\), the operations
\[
 K_j(t,s)=U(t)\ee^{-\gamma(t-s)X^2/2}
       \frac{(\sqrt{\gamma(t-s)}X)^j}{\sqrt{j!}}U(s)^*
\]
are complete since the positive sum \(\sum_jK_j^*K_j\) is the identity
by the scalar exponential series. The no-count amplitude multiplies
\(f_{1/2}\) by \(\ee^{-\gamma t u^2/2}\), giving
\(\lambda=1/2-\gamma t\), \(0\le t\le1/\gamma\). The effects
\((I\pm\cos(kQ_t))/2\) then obey
\[
 q_++q_-=H_\lambda(0)/H_{1/2}(0),\qquad
 H_\lambda(k)=H_{1/2}(0)(q_+-q_-).
\]
These are joint probabilities including failed conditioning; arbitrary
complex \(k\) is not claimed to define a bounded effect.

In the auxiliary model the supplied amplitude is instead
\(W_\lambda(v)B\), where \(BB^*=W\ee^{-\pi A}/Q(\pi)\) and
\(\|B\|_{\rm HS}=1\). Apply \(\ee^{-\tau L_0}\) on the physical side
and pair the Fourier amplitude with \(B\). Commuting radial and angular
heat factors and summing the complete angular trace gives
\begin{equation}\label{eq:auxiliary-response}
 \mathcal O_{\lambda,\tau}(z)
 =\ee^{-\tau z^2}\frac{Q(\pi+4\tau)}{Q(\pi)}H_\lambda(z).
\end{equation}
The multiplier never vanishes, so this response preserves zeros and
multiplicities. Recovery on real \(z=k\) amplifies errors by
\(\ee^{\tau k^2}Q(\pi)/Q(\pi+4\tau)\), unbounded in \(k\) for
\(\tau>0\). Proper time in this model is not Newman time.
The graph-domain and weighted-integral proofs for both models remain
in the parent sources; their loaders, quadrature access and calibrated
detector laws are supplied assumptions. Their exact observation formulas
do not generate a missing source-specific inequality.
